\section{Introduction}
Earth observation has become an indispensable tool for understanding and monitoring the planet's dynamic processes. Optical and radar remote sensing represent two complementary modalities at the core of modern monitoring systems. Optical sensors, such as the Multi-Spectral Instrument (MSI) onboard Sentinel-2, provide rich spectral and visual information essential for applications ranging from agricultural monitoring and land-cover classification to disaster response and forest monitoring~\cite{S2MS_GAN}. However, these sensors are inherently constrained by atmospheric conditions, particularly cloud cover. Global estimates suggest that average cloud cover exceeds 66\%~\cite{dl_cloud_detection_survey,A_cGAN_fuse_sar_MS_CR,CR_SEN2_dRNN}, with about 55\% over land surfaces~\cite{CR_SEN2_dRNN}, causing considerable data gaps in both spatial and temporal domains. For applications requiring consistent time series, e.g., agricultural monitoring, or where a specific scene must be observed at a given time, e.g., disaster monitoring, cloud cover represents a serious limitation~\cite{CR_SEN2_dRNN}. The diversity of clouds —including thin and thick clouds as well as haze— together with the wide range of occlusion scenarios and their uneven distribution, poses an additional challenge for image reconstruction and the generalizability of cloud removal techniques~\cite{CR_Advances_Review_ORS}. In contrast, Synthetic Aperture Radar (SAR) sensors, operating in the microwave domain, such as the C-band instrument on Sentinel-1 mission offer all-weather, day-and-night imaging capabilities independent of solar illumination or cloud interference. However, the backscatter-based nature of SAR imagery introduces challenges related to speckle noise, geometric distortions, and a lack of intuitive spectral color information \cite{bench_sar_color, CR_RS_GAN_s2o}, necessitating advanced interpretation skills.

To bridge the gap between these two sensing modalities, SAR-to-optical image translation has emerged as a powerful generative approach. It aims to synthesize optical-like, cloud-free imagery from SAR data, combining the interpretability of optical observations with the reliability of radar acquisitions. The domain gap between SAR and optical imagery, however, poses significant challenges. For example, while two objects with identical structures may appear different in optical imagery due to their spectral responses, they can appear similar in SAR imagery, reflecting SAR’s emphasis on structural rather than spectral properties~\cite{RS_Data_Fusion_GANs_sota}. Moreover, acquiring perfectly co-registered SAR–optical pairs is also non-trivial, as both spatial and temporal alignment must be ensured. 

Recent advances in generative artificial intelligence (GenAI), particularly in conditional generative adversarial networks (cGANs)~\cite{cGANs_mirza} and diffusion models~\cite{DDPM_2020}, have made it possible to learn complex mappings between SAR and optical domains with visually plausible results~\cite{cr_fuse_HR_GEN}. These developments have opened new pathways for applications in land-cover classification, vegetation monitoring, disaster response, and particularly, cloud removal. Despite the rapid progress in this field, a significant gap remains in the literature. Most existing studies have focused primarily on reconstructing the visible RGB subset of optical imagery such as~\cite{CR_RS_GAN_s2o, cr_fuse_HR_GEN, c_diffusion_s2o, hvt_cgan, s2o_color_super_diff}, with a few extending to the Near-Infrared (NIR) range such as~\cite{filmy_cloudy_NIR_Enomoto, Mult_temp_S12_fus}. Consequently, the full multispectral potential of missions such as Sentinel-2 remains largely underexplored in the context of generative reconstruction. Critical bands for vegetation analysis (Red-Edge) and soil moisture monitoring (SWIR) are rarely addressed. 

Furthermore, while SAR-to-optical translation is often evaluated qualitatively, systematic assessments of its reliability across individual spectral bands remain scarce. Finally, although the approach shows promise for cloud removal, its effectiveness in this context has not yet been comprehensively validated across the full spectral range.

Accordingly, this study investigates the use of generative models for translating Sentinel-1 SAR imagery into the full multispectral stack of Sentinel-2 imagery. The objectives of this work are threefold:
\begin{enumerate}
    \item  To validate SAR-to-optical image translation across all 13 Sentinel-2 spectral bands, moving beyond standard RGB constraints.
    \item To systematically assess the reliability and reconstructability of each individual spectral band.
    \item  To evaluate the performance of the translation model for the practical application of cloud removal.
\end{enumerate}
By addressing these objectives, this work aims to provide a comprehensive and quantitative understanding of SAR-to-optical translation as a multimodal learning problem and to clarify its potential and limitations for enhancing the temporal continuity of Earth observation data.